\documentclass[12pt]{article}
\usepackage[utf8]{inputenc}
\usepackage[T1]{fontenc}
\usepackage[margin=1in]{geometry}
\usepackage{setspace}
\usepackage{xcolor}
\usepackage[colorlinks=true, urlcolor=blue!60!black, linkcolor=blue!60!black]{hyperref}
\setlength{\parindent}{2em}
\setlength{\parskip}{1.25em}
\renewcommand{\baselinestretch}{1.0}
\begin{document}
\begin{center}
{\LARGE \textit{Garcetti v. Ceballos}, 547 U.S. 410 (2006)\par}
\vspace{0.8em}
{\large Case Brief\par}
\end{center}
\vspace{1em}

\section{Memory Jogger}

When a public employee speaks pursuant to official job duties---such as a deputy district attorney writing an internal case-disposition memorandum questioning a search warrant affidavit---the employee is not speaking as a private citizen, and the First Amendment provides no protection against employer discipline (\href{https://www.courtlistener.com/opinion/145653/garcetti-v-ceballos/}{CourtListener}; \href{https://supreme.justia.com/cases/federal/us/547/410/}{Justia}).

\section{Detailed Case Facts}

Richard Ceballos worked as a deputy district attorney for the Los Angeles County District Attorney's Office beginning in 1989 (\href{https://supreme.justia.com/cases/federal/us/547/410/}{Justia}). During the events at issue, he served as a calendar deputy in the office's Pomona branch, a role carrying supervisory responsibilities over other attorneys (\href{https://supreme.justia.com/cases/federal/us/547/410/}{Justia}). In February 2000, a defense attorney contacted Ceballos about a pending criminal prosecution, claiming that the affidavit police used to obtain a critical search warrant contained inaccuracies and notifying him that the defense had filed a motion to traverse (challenge) the warrant (\href{https://supreme.justia.com/cases/federal/us/547/410/}{Justia}).

After examining the affidavit and visiting the location it described, Ceballos concluded that the affidavit contained serious misrepresentations: it described as a ``long driveway'' what Ceballos considered a separate roadway, and it asserted that tire tracks led from a stripped-down truck to the premises even though the roadway's composition made visible tire tracks difficult or impossible in places (\href{https://supreme.justia.com/cases/federal/us/547/410/}{Justia}). Ceballos telephoned the warrant affiant, a deputy sheriff with the Los Angeles County Sheriff's Department, but received no satisfactory explanation (\href{https://supreme.justia.com/cases/federal/us/547/410/}{Justia}).

Ceballos reported his findings to his supervisors, Carol Najera and Head Deputy District Attorney Frank Sundstedt, and on March 2, 2000, submitted a disposition memorandum to Sundstedt recommending dismissal of the criminal case (\href{https://supreme.justia.com/cases/federal/us/547/410/}{Justia}). A few days later, he submitted a second memorandum recounting another phone call with the affiant (\href{https://supreme.justia.com/cases/federal/us/547/410/}{Justia}). A heated meeting followed among Ceballos, Sundstedt, Najera, the warrant affiant, and other Sheriff's Department officials, during which a sheriff's lieutenant sharply criticized Ceballos's handling of the matter (\href{https://supreme.justia.com/cases/federal/us/547/410/}{Justia}). Sundstedt nonetheless decided to proceed with the prosecution pending resolution of the defense motion (\href{https://supreme.justia.com/cases/federal/us/547/410/}{Justia}). At the suppression hearing, the defense called Ceballos to recount his observations about the affidavit, though the trial court denied the motion and upheld the warrant (\href{https://supreme.justia.com/cases/federal/us/547/410/}{Justia}).

Ceballos alleged that his supervisors---including District Attorney Gil Garcetti, Sundstedt, and Najera---subsequently subjected him to a series of retaliatory employment actions: reassigning him from calendar deputy to trial deputy, transferring him to another courthouse (the El Monte branch), and denying him a promotion (\href{https://supreme.justia.com/cases/federal/us/547/410/}{Justia}). After his internal employment grievance was denied on the ground that he had suffered no retaliation, Ceballos sued under 42 U.S.C.\ \S~1983, alleging violations of the First and Fourteenth Amendments (\href{https://supreme.justia.com/cases/federal/us/547/410/}{Justia}).

\section{Procedural History}

\begin{enumerate}
  \item \textbf{U.S.\ District Court for the Central District of California:} Ceballos filed a \S~1983 action alleging unconstitutional retaliation for his March 2, 2000 disposition memorandum (\href{https://supreme.justia.com/cases/federal/us/547/410/}{Justia}). Petitioners denied retaliation, attributing the personnel changes to legitimate staffing needs, and moved for summary judgment (\href{https://supreme.justia.com/cases/federal/us/547/410/}{Justia}). On January 30, 2002, the District Court granted summary judgment to petitioners, holding that because Ceballos wrote the memorandum pursuant to his employment duties, it was unprotected by the First Amendment, and holding in the alternative that petitioners were entitled to qualified immunity (\href{https://en.wikipedia.org/wiki/Garcetti_v._Ceballos}{Wikipedia}; \href{https://supreme.justia.com/cases/federal/us/547/410/}{Justia}).
  \item \textbf{U.S.\ Court of Appeals for the Ninth Circuit:} In \textit{Ceballos v. Garcetti}, 361 F.3d 1168 (9th Cir.\ 2004), a panel reversed (\href{https://supreme.justia.com/cases/federal/us/547/410/}{Justia}). Applying \textit{Pickering v. Board of Education}, 391 U.S. 563 (1968), and \textit{Connick v. Myers}, 461 U.S. 138 (1983), the Ninth Circuit held that Ceballos's allegations of government misconduct addressed a matter of public concern, declined to exclude speech made pursuant to job duties under circuit precedent, struck the \textit{Pickering} balance in Ceballos's favor, and denied qualified immunity (\href{https://supreme.justia.com/cases/federal/us/547/410/}{Justia}). Judge Diarmuid O'Scannlain specially concurred, bound by circuit precedent but urging that routine job-duty speech should receive no First Amendment protection (\href{https://supreme.justia.com/cases/federal/us/547/410/}{Justia}).
  \item \textbf{U.S.\ Supreme Court:} The Court granted certiorari on February 28, 2005 (543 U.S. 1186), heard initial argument on October 12, 2005, and ordered reargument on March 21, 2006 following Justice Sandra Day O'Connor's retirement and Justice Samuel A. Alito, Jr.'s confirmation (\href{https://supreme.justia.com/cases/federal/us/547/410/}{Justia}; \href{https://open.mitchellhamline.edu/cgi/viewcontent.cgi?article=1237&context=mhlr}{Mitchell Hamline L. Rev.}). On May 30, 2006, the Supreme Court reversed and remanded (\href{https://supreme.justia.com/cases/federal/us/547/410/}{Justia}).
\end{enumerate}

\section{Judicial Votes}

\textbf{Vote Breakdown:} 5--4, reversing and remanding the Ninth Circuit (\href{https://supreme.justia.com/cases/federal/us/547/410/}{Justia}; \href{https://tile.loc.gov/storage-services/service/ll/usrep/usrep547/usrep547410/usrep547410.pdf}{Library of Congress U.S.\ Reports}).

\begin{itemize}
  \item \textbf{Majority (5):} Justice Anthony M. Kennedy (author), joined by Chief Justice John G. Roberts, Jr., and Justices Antonin Scalia, Clarence Thomas, and Samuel A. Alito, Jr.
  \item \textbf{Concurrences (0):} None.
  \item \textbf{Dissents (4):}
  \begin{itemize}
    \item Justice John Paul Stevens (filed a separate dissenting opinion; also joined Justice Souter's dissent).
    \item Justice David H. Souter (filed a dissenting opinion, joined by Justices Stevens and Ruth Bader Ginsburg).
    \item Justice Stephen G. Breyer (filed a separate dissenting opinion).
  \end{itemize}
\end{itemize}

\section{Holding}

\textbf{Judgment:} Reversed and remanded (\href{https://supreme.justia.com/cases/federal/us/547/410/}{Justia}).

\textbf{Rule Announced:} When public employees make statements pursuant to their official duties, they are not speaking as citizens for First Amendment purposes, and the Constitution does not insulate their communications from employer discipline (\href{https://supreme.justia.com/cases/federal/us/547/410/}{Justia}). This ``official duties'' threshold precedes \textit{Pickering} balancing: only if an employee speaks (1) as a citizen (2) on a matter of public concern does a court proceed to ask whether the government employer had an adequate justification for treating the employee differently from any other member of the general public (\href{https://supreme.justia.com/cases/federal/us/547/410/}{Justia}).

\textbf{Holding vs.\ Dicta:} The narrow holding resolves only undisputed official-duty speech: because the parties agreed that Ceballos prepared his disposition memorandum pursuant to his duties as a calendar deputy, the memorandum lacked First Amendment protection (\href{https://supreme.justia.com/cases/federal/us/547/410/}{Justia}). Because the scope of duties was undisputed, the Court expressly declined to articulate a comprehensive test for defining official duties, though it stated in influential guidance that the inquiry is ``a practical one'' and that formal written job descriptions are ``neither necessary nor sufficient'' (\href{https://supreme.justia.com/cases/federal/us/547/410/}{Justia}). The Court also expressly reserved---and thus left undecided as a reservation rather than a holding---whether its official-duties rule applies in the same manner to ``speech related to scholarship or teaching'' in public colleges and universities (\href{https://supreme.justia.com/cases/federal/us/547/410/}{Justia}). Finally, the Court's observations that statutory whistleblower protections, labor codes, and professional ethics rules adequately check governmental misconduct serve as policy support rather than binding holdings (\href{https://supreme.justia.com/cases/federal/us/547/410/}{Justia}).

\section{Analysis of Opinions}

\subsection{Majority Opinion (Justice Kennedy)}

Justice Kennedy framed the case around the two-step \textit{Pickering}--\textit{Connick} inquiry: first, whether the employee spoke as a citizen on a matter of public concern; second, if so, whether the government employer had an adequate justification for restricting the speech (\href{https://supreme.justia.com/cases/federal/us/547/410/}{Justia}). Emphasizing that neither the internal location of Ceballos's speech (\textit{Givhan v. Western Line Consolidated School District}, 439 U.S. 410 (1979)) nor its job-related subject matter (\textit{Pickering}, 391 U.S. at 572) was dispositive, Justice Kennedy identified the controlling factor as whether the expression was made pursuant to Ceballos's official duties (\href{https://supreme.justia.com/cases/federal/us/547/410/}{Justia}).

Drawing an analogy to government speech doctrine in \textit{Rosenberger v. Rector and Visitors of the University of Virginia}, 515 U.S. 819, 833 (1995), the majority reasoned that restricting speech that ``owes its existence to a public employee's professional responsibilities'' does not abridge private citizen liberties; rather, it reflects ``employer control over what the employer itself has commissioned or created'' (\href{https://supreme.justia.com/cases/federal/us/547/410/}{Justia}). Official communications have official consequences, requiring substantive consistency, accuracy, and managerial oversight (\href{https://supreme.justia.com/cases/federal/us/547/410/}{Justia}). Extending First Amendment claims to official-duty speech would invite permanent judicial supervision of internal agency operations, offending federalism and separation of powers (\href{https://supreme.justia.com/cases/federal/us/547/410/}{Justia}). The majority rejected the Ninth Circuit's concern about a ``doctrinal anomaly'' (protecting public leaks while leaving internal complaints unprotected), noting that public statements have a citizen analogue whereas official work product does not, and that employers can adopt internal complaint channels (\href{https://supreme.justia.com/cases/federal/us/547/410/}{Justia}). Finally, the majority pointed to federal and state whistleblower laws, labor codes, attorney ethics rules, and constitutional duties under \textit{Brady v. Maryland}, 373 U.S. 83 (1963), as alternative safeguards against misconduct (\href{https://supreme.justia.com/cases/federal/us/547/410/}{Justia}).

\subsection{Dissenting Opinion of Justice Stevens}

In a brief, pointed dissent, Justice Stevens argued that the proper answer to whether the First Amendment protects official-duty speech is ``Sometimes,'' not ``Never'' (\href{https://supreme.justia.com/cases/federal/us/547/410/}{Justia}). While supervisors may discipline ``inflammatory or misguided'' work product, they should not be free to punish unwelcome speech solely because it uncovers facts a supervisor wishes to conceal (\href{https://supreme.justia.com/cases/federal/us/547/410/}{Justia}). Citing \textit{Givhan}, Justice Stevens contended that public employees remain citizens inside the workplace and that it is ``senseless'' to make First Amendment protection for identical words turn on whether they fall within a job description---a rule that perversely incentivizes employees to go straight to the press rather than speaking candidly with superiors (\href{https://supreme.justia.com/cases/federal/us/547/410/}{Justia}).

\subsection{Dissenting Opinion of Justice Souter (Joined by Justices Stevens and Ginsburg)}

Justice Souter argued that the majority's categorical exclusion sacrificed both individual and public interests in exposing official wrongdoing and threats to health and safety (\href{https://supreme.justia.com/cases/federal/us/547/410/}{Justia}). Rather than a ``winner-take-all'' bar, Justice Souter proposed an adjusted \textit{Pickering} balance under which official-duty speech would prevail only when the employee speaks on a matter of ``unusual importance''---such as official dishonesty, deliberately unconstitutional action, serious wrongdoing, or threats to health and safety---and satisfies high standards of responsibility (\href{https://supreme.justia.com/cases/federal/us/547/410/}{Justia}).

Justice Souter dismantled two pillars of the majority's reasoning. First, he attacked the majority's reliance on \textit{Rust v. Sullivan}, 500 U.S. 173 (1991), and \textit{Rosenberger}, explaining that Ceballos was not hired to broadcast a preset government manifesto, but to enforce the law honestly and constitutionally; equating all official work product with government speech risked ``a bloated notion of controllable government speech'' that could imperil academic freedom at public universities (\href{https://supreme.justia.com/cases/federal/us/547/410/}{Justia}). Second, he demonstrated that statutory whistleblower protections form an inconsistent ``patchwork''---for example, federal courts had interpreted the federal Whistleblower Protection Act to exclude disclosures made to immediate supervisors or in the normal course of employment duties (\href{https://supreme.justia.com/cases/federal/us/547/410/}{Justia}). He also warned that employers might expand written job descriptions to swallow citizen speech, and noted that on remand the Ninth Circuit remained free to evaluate Ceballos's other speech, including his courtroom testimony and his speech to the Mexican-American Bar Association (\href{https://supreme.justia.com/cases/federal/us/547/410/}{Justia}).

\subsection{Dissenting Opinion of Justice Breyer}

Justice Breyer rejected both the majority's categorical ``never'' and Justice Souter's modified \textit{Pickering} test (\href{https://supreme.justia.com/cases/federal/us/547/410/}{Justia}). He feared that Justice Souter's ``unusual importance'' standard would sweep in too much routine administrative speech by police officers, inspectors, and regulators, inviting discovery and judicial interference with agency management (\href{https://supreme.justia.com/cases/federal/us/547/410/}{Justia}). Instead, Justice Breyer advocated applying \textit{Pickering} balancing to official-duty speech only where two special conditions converge: (1) heightened need for constitutional protection and (2) diminished risk of undue judicial interference (\href{https://supreme.justia.com/cases/federal/us/547/410/}{Justia}). Ceballos met that narrow test because his speech was governed by independent professional canons for lawyers and by an affirmative constitutional duty under \textit{Brady v. Maryland} to disclose exculpatory evidence (\href{https://supreme.justia.com/cases/federal/us/547/410/}{Justia}).

\section{Examples: Future Applications}

\begin{enumerate}
  \item \textbf{Same-Side (Unprotected Official-Duty Speech):} A municipal building inspector submits a required internal inspection report warning her supervisor that a politically connected developer's apartment complex violates fire codes, and is demoted for causing political embarrassment. Under \textit{Garcetti}, because inspecting buildings and filing compliance reports are the exact tasks she is paid to perform, her report owes its existence to her official duties and receives no First Amendment protection---leaving her to rely solely on state whistleblower statutes (\href{https://supreme.justia.com/cases/federal/us/547/410/}{Justia}).
  \item \textbf{Same-Side (Unprotected Official-Duty Speech):} A public high school teacher departs from the school board's approved curriculum during an assigned history class to deliver a personal political critique of current foreign policy and is suspended. Because classroom instruction in a K--12 public school is speech delivered pursuant to the teacher's official duties to transmit the state's curriculum, lower courts applying \textit{Garcetti} generally treat K--12 instructional speech as unprotected official-duty speech (\href{https://onlabor.org/garcettis-impact-on-teachers/}{OnLabor}; \href{https://digitalcommons.pepperdine.edu/cgi/viewcontent.cgi?article=2685&context=plr}{Pepperdine L. Rev.}).
  \item \textbf{Opposite-Side (Protected Citizen Speech):} A state community college program director uncovers payroll fraud during an audit and is later subpoenaed to give truthful sworn testimony at a federal criminal trial against the fraudster; the college president fires him in retaliation. As the Supreme Court later held in \textit{Lane v. Franks}, 573 U.S. 228 (2014), truthful sworn testimony outside an employee's ordinary job responsibilities is citizen speech on a matter of public concern---even when it concerns information learned on the job---and therefore triggers \textit{Pickering} protection (\href{https://www.journaloffreespeechlaw.org/kozel.pdf}{Journal of Free Speech Law}; \href{https://digitalcommons.pepperdine.edu/cgi/viewcontent.cgi?article=2685&context=plr}{Pepperdine L. Rev.}).
  \item \textbf{Opposite-Side (Protected Citizen Speech):} A public high school football coach kneels for a brief, quiet personal prayer at midfield after games during a post-game period when coaching staff are free to attend to personal matters, and is disciplined by the school district. As the Supreme Court held in \textit{Kennedy v. Bremerton School District}, 597 U.S. 507 (2022), because the coach was not conveying a government-created message or speaking pursuant to his official coaching duties during that personal window, \textit{Garcetti} does not bar his First Amendment claim (\href{https://uknowledge.uky.edu/cgi/viewcontent.cgi?article=5590&context=klj}{Kentucky L.J.}; \href{https://www.journaloffreespeechlaw.org/kozel.pdf}{Journal of Free Speech Law}).
  \item \textbf{Fence-Sitter (Borderline Case):} A tenured chemistry professor at a public state university publishes a peer-reviewed article criticizing the state environmental agency's water-safety standards and discusses the findings in her seminar, prompting university trustees to deny her a merit raise. Under \textit{Garcetti}, scholarship and classroom teaching are plainly part of what a public professor is paid to do, which would ordinarily foreclose First Amendment review; yet Justice Kennedy explicitly reserved whether \textit{Garcetti} applies ``in the same manner to a case involving speech related to scholarship or teaching'' (\href{https://supreme.justia.com/cases/federal/us/547/410/}{Justia}; \href{https://fedsoc.org/fedsoc-review/an-academic-freedom-exception-to-government-control-of-employee-speech}{Federalist Society Review}). Whether the professor wins turns on whether the reviewing circuit recognizes an academic-freedom exception to \textit{Garcetti}.
\end{enumerate}

\section{Critique}

\textbf{Originalist and Textualist Perspectives:} Supporters of \textit{Garcetti} rooted in constitutional structure and formalism emphasize that the First Amendment's text (``Congress shall make no law \ldots\ abridging the freedom of speech'') restrains the government acting as \textit{sovereign}, not as an \textit{employer} managing its own workforce, and that until the mid-twentieth century the traditional understanding permitted broad governmental control over public employment (\href{https://scholarship.law.nd.edu/context/ndlr/article/1016/viewcontent/21_87NotreDameLRev841_2011_2012_.pdf}{Notre Dame L. Rev.}; \href{https://fedsoc.org/case/garcetti-v-ceballos}{Federalist Society}). On this view, \textit{Garcetti}'s bright-line rule preserves separation of powers and federalism by preventing federal judges from micromanaging state and local executive agencies (\href{https://scholarship.law.nd.edu/context/ndlr/article/1016/viewcontent/21_87NotreDameLRev841_2011_2012_.pdf}{Notre Dame L. Rev.}). Yet other formalist and textualist critics argue that \textit{Garcetti} overextends government-speech doctrine without textual grounding: while the government must control messages it hires spokespersons to deliver, equating all on-duty expression with ``government speech'' ignores the First Amendment's text protecting speakers and listeners alike, and its vague ``practical inquiry'' has ironically generated deep circuit splits rather than formalist predictability (\href{https://scholarship.law.nd.edu/context/ndlr/article/1016/viewcontent/21_87NotreDameLRev841_2011_2012_.pdf}{Notre Dame L. Rev.}; \href{https://fedsoc.org/fedsoc-review/an-academic-freedom-exception-to-government-control-of-employee-speech}{Federalist Society Review}).

\textbf{Progressive and Doctrinal Perspectives:} Progressive and First Amendment scholars have sharply criticized \textit{Garcetti} as a severe contraction of public-sector speech rights (\href{https://open.mitchellhamline.edu/cgi/viewcontent.cgi?article=1237&context=mhlr}{Mitchell Hamline L. Rev.}). Professor Helen Norton contends that \textit{Garcetti} rests on a distorted theory of government speech, allowing the state to claim employees' on-duty speech as its own without transparency or political accountability, thereby depriving the public of vital insider information about agency misconduct (\href{https://scholar.law.colorado.edu/context/faculty-articles/article/1915/viewcontent/2008_Norton_Government_Workers_and_Government_Speech.pdf}{Colorado Law}; \href{https://www.journaloffreespeechlaw.org/kozel.pdf}{Journal of Free Speech Law}). Doctrinal critics also highlight the perverse incentive noted by Justice Stevens: employees who report wrongdoing responsibly up the chain of command lose constitutional protection, whereas employees who leak to the media may retain \textit{Pickering} protection (\href{https://supreme.justia.com/cases/federal/us/547/410/}{Justia}; \href{https://scholarship.law.unc.edu/cgi/viewcontent.cgi?article=1087&context=falr}{First Amendment L. Rev.}). Moreover, scholars document that statutory whistleblower remedies are riddled with gaps and procedural hurdles, leaving conscientious public servants---and potentially public university faculty---vulnerable to retaliation (\href{https://scholarship.law.unc.edu/cgi/viewcontent.cgi?article=1087&context=falr}{First Amendment L. Rev.}; \href{https://digitalcommons.law.seattleu.edu/cgi/viewcontent.cgi?article=2199&context=sulr}{Seattle U. L. Rev.}).

\textbf{Underlying Competing Values:} \textit{Garcetti} exposes a fundamental tension between \textbf{democratic accountability through managerial control} (elected officials and agency supervisors must be able to direct employees' work product to execute public policy coherently and efficiently) and \textbf{democratic accountability through transparency} (the public depends on knowledgeable government insiders to expose illegality, corruption, and safety threats without fear of reprisal).

\section{Key Quotations}

\begin{quote}
``We hold that when public employees make statements pursuant to their official duties, the employees are not speaking as citizens for First Amendment purposes, and the Constitution does not insulate their communications from employer discipline.''
\end{quote}
\noindent --- Justice Anthony M. Kennedy, Opinion of the Court, 547 U.S. at 421 (\href{https://supreme.justia.com/cases/federal/us/547/410/}{Justia}; \href{https://www.courtlistener.com/opinion/145653/garcetti-v-ceballos/}{CourtListener}).

\begin{quote}
``Restricting speech that owes its existence to a public employee's professional responsibilities does not infringe any liberties the employee might have enjoyed as a private citizen. It simply reflects the exercise of employer control over what the employer itself has commissioned or created.''
\end{quote}
\noindent --- Justice Anthony M. Kennedy, Opinion of the Court, 547 U.S. at 421--422 (\href{https://supreme.justia.com/cases/federal/us/547/410/}{Justia}; \href{https://tile.loc.gov/storage-services/service/ll/usrep/usrep547/usrep547410/usrep547410.pdf}{Library of Congress U.S.\ Reports}).

\begin{quote}
``The proper inquiry is a practical one. Formal job descriptions often bear little resemblance to the duties an employee actually is expected to perform, and the listing of a given task in an employee's written job description is neither necessary nor sufficient to demonstrate that conducting the task is within the scope of the employee's professional duties for First Amendment purposes.''
\end{quote}
\noindent --- Justice Anthony M. Kennedy, Opinion of the Court, 547 U.S. at 424--425 (\href{https://supreme.justia.com/cases/federal/us/547/410/}{Justia}).

\begin{quote}
``As Justice Souter explains, public employees are still citizens while they are in the office. The notion that there is a categorical difference between speaking as a citizen and speaking in the course of one's employment is quite wrong.''
\end{quote}
\noindent --- Justice John Paul Stevens, dissenting (\href{https://supreme.justia.com/cases/federal/us/547/410/}{Justia}).

\begin{quote}
``Even there, the lesson of \textit{Pickering} (and the object of most constitutional adjudication) is still to the point: when constitutionally significant interests clash, resist the demand for winner-take-all; try to make adjustments that serve all of the values at stake.''
\end{quote}
\noindent --- Justice David H. Souter, dissenting, joined by Justices Stevens and Ginsburg (\href{https://supreme.justia.com/cases/federal/us/547/410/}{Justia}).

\end{document}
